\documentclass[12pt,a4paper]{article}
\usepackage[utf8]{inputenc}
\usepackage[T1]{fontenc}
\usepackage[polish]{babel}
\usepackage{geometry}
\geometry{a4paper, margin=2.5cm}
\usepackage{graphicx}
\usepackage{amsmath}
\usepackage{booktabs}
\usepackage{listings}
\usepackage{xcolor}
\usepackage{hyperref}

\lstset{
    basicstyle=\ttfamily\footnotesize,
    backgroundcolor=\color{gray!10},
    frame=single,
    breaklines=true,
    showstringspaces=false,
    language=Python,
    commentstyle=\color{green!50!black},
    keywordstyle=\color{blue},
    stringstyle=\color{red!70!black},
}

\begin{document}

% ===================== STRONA TYTUŁOWA =====================
\thispagestyle{empty}
\begin{center}
\begin{tabular}{|p{5cm}|p{9cm}|}
\hline
\textbf{Uczelnia} & Politechnika Bydgoska im. J. J. Śniadeckich \\ \hline
\textbf{Wydział} & Telekomunikacji, Informatyki i Elektrotechniki \\ \hline
\textbf{Przedmiot} & Sztuczna Inteligencja \\ \hline
\textbf{Tytuł projektu} & System AI do gry w Kuhn Poker oparty na sieci neuronowej i algorytmie genetycznym \\ \hline
\textbf{Autorzy} & [Imię i Nazwisko 1] \newline [Imię i Nazwisko 2] \\ \hline
\textbf{Grupa} & 4 \\ \hline
\textbf{Data} & \today \\ \hline
\textbf{Prowadzący} & dr inż. Tomasz Talaśka \\ \hline
\end{tabular}
\end{center}
\vfill
\begin{center}
\Large Rok akademicki 2025/2026
\end{center}
\newpage

% ===================== SEKCJA 1 =====================
\section{Cel projektu i ogólny opis}

Celem projektu jest zaprojektowanie i implementacja systemu sztucznej
inteligencji uczącego się grać w \emph{Kuhn Poker} -- uproszczoną,
dwuosobową odmianę pokera, powszechnie wykorzystywaną w literaturze
teorii gier jako modelowy przykład gry z niepełną informacją. Sercem
rozwiązania jest sieć neuronowa pełniąca funkcję polityki gracza, której
wagi są optymalizowane algorytmem genetycznym (neuroewolucja). Takie
podejście pozwala uniknąć konieczności zbierania zbioru uczącego oraz
różniczkowania funkcji nagrody, co jest szczególnie istotne w grach
sekwencyjnych z elementem losowości.

Projekt łączy dwa zagadnienia z programu przedmiotu: sieci neuronowe
(temat 3) oraz algorytmy ewolucyjne wraz z systemami ekspertowymi
(temat 4). Jako punkt odniesienia zaimplementowany został również prosty
system ekspertowy oraz strategia Nasha o znanej postaci analitycznej,
co umożliwia jakościową i ilościową ocenę uzyskanego agenta.

\subsection*{Założenia projektowe}
\begin{itemize}
    \item Wariant gry: Kuhn Poker, dwóch graczy, talia trzech kart (J, Q, K).
    \item Metoda główna: sieć neuronowa typu MLP + algorytm genetyczny.
    \item Punkty odniesienia: agent losowy, system ekspertowy, strategia Nasha.
    \item Miary jakości: średnia wygrana na rozdanie, krzywa zbieżności
          fitness, wynik bezpośredniej konfrontacji ze strategiami referencyjnymi.
    \item Język i biblioteki: Python, NumPy, DEAP (algorytm genetyczny), PyTorch (sieć neuronowa), Matplotlib (wizualizacja wyników).
\end{itemize}

\subsection*{Spodziewany wynik}
Oczekujemy, że agent po treningu uzyska wynik zbliżony do strategii Nasha
w konfrontacji ze słabszymi przeciwnikami (losowy, zawsze-bet, system
ekspertowy o ograniczonych regułach), a w grze przeciwko optymalnej
strategii Nasha jego strata na rozdanie będzie bliska teoretycznej
wartości wyjściowej gry, czyli okolicy $-1/18 \approx -0.0556$ żetona
na rozdanie (z perspektywy gracza działającego pierwszego).

% ===================== SEKCJA 2 =====================
\section{Dane wejściowe -- opis problemu i wyjaśnienie algorytmu}

\subsection{Reguły Kuhn Poker}
Kuhn Poker jest zredukowaną wersją pokera klasycznego, zaproponowaną
przez Harolda W.~Kuhna w 1950 r. Reguły:
\begin{enumerate}
    \item Talia składa się z trzech kart: Jack (J=1), Queen (Q=2), King (K=3).
    \item Każdy z dwóch graczy wnosi do puli 1 żeton (\emph{ante}).
    \item Każdemu graczowi rozdawana jest jedna zakryta karta; trzecia karta pozostaje w talii.
    \item Gracz~1 wykonuje pierwszy ruch. Dostępne akcje:
          \emph{pass} (\texttt{p}) oraz \emph{bet} (\texttt{b}).
          Znaczenie zależy od kontekstu: \texttt{p} to \emph{check} w sytuacji bez zakładu i \emph{fold} po zakładzie przeciwnika; \texttt{b} to \emph{bet} w sytuacji bez zakładu i \emph{call} po zakładzie przeciwnika.
    \item Gra kończy się jednym z poniższych terminalnych ciągów akcji:
\end{enumerate}

\begin{center}
\begin{tabular}{lll}
\toprule
Historia & Opis & Wypłata gracza 1 \\ \midrule
\texttt{pp}  & obaj sprawdzają        & showdown, $\pm 1$ \\
\texttt{pbp} & P1 sprawdza, P2 bet, P1 fold & $-1$ \\
\texttt{pbb} & P1 sprawdza, P2 bet, P1 call & showdown, $\pm 2$ \\
\texttt{bp}  & P1 bet, P2 fold        & $+1$ \\
\texttt{bb}  & P1 bet, P2 call        & showdown, $\pm 2$ \\
\bottomrule
\end{tabular}
\end{center}

W \emph{showdown} wygrywa gracz posiadający wyższą kartę.

\subsection{Reprezentacja danych wejściowych}
Stan gry widoczny dla gracza składa się z dwóch elementów: jego własnej
karty oraz dotychczasowej historii akcji. Kodujemy go w wektorze o
długości 12:
\begin{itemize}
    \item Wymiary $[0,3)$: \emph{one-hot} dla karty (J, Q, K).
    \item Wymiary $[3,6)$, $[6,9)$, $[9,12)$: \emph{one-hot} dla każdej
          z trzech możliwych pozycji w historii akcji
          (kategorie: brak ruchu, pass, bet).
\end{itemize}
Maksymalna długość historii przed terminalem to 3 akcje, zatem
dwunastowymiarowy wektor pokrywa wszystkie sytuacje decyzyjne.

\subsection{Algorytm: neuroewolucja}
Algorytm uczenia łączy dwa komponenty:

\paragraph{1. Polityka -- sieć neuronowa.}
Sieć perceptronowa typu MLP przyjmuje na wejściu wektor stanu (12
wymiarów), a na wyjściu generuje 2 logity odpowiadające akcjom
\emph{pass} i \emph{bet}. Akcja wybierana jest jako $\arg\max$ z logitów.
Wagi sieci są spłaszczane do pojedynczego wektora rzeczywistego, który
jest genomem osobnika populacji.

\paragraph{2. Optymalizacja -- algorytm genetyczny.}
Korzystamy z biblioteki DEAP. Schemat działania:
\begin{enumerate}
    \item Inicjalizacja populacji $P$ losowych wektorów wag.
    \item Dla każdego osobnika obliczana jest funkcja fitness: agent
          gra ustaloną liczbę rozdań przeciwko puli przeciwników
          (system ekspertowy, agent losowy, agent zawsze-bet),
          a wynik to średnia wygrana na rozdanie. Pozycja przy stole
          jest naprzemienna, by wyeliminować przewagę pierwszego ruchu.
    \item Selekcja turniejowa o rozmiarze $k=3$.
    \item Krzyżowanie \emph{blend} ($\alpha=0{,}5$).
    \item Mutacja gaussowska ($\mu=0$, $\sigma=0{,}2$, prawdopodobieństwo per-gen $=0{,}1$).
    \item Elityzm: najlepszy osobnik (Hall of Fame) zachowywany w każdej iteracji.
    \item Powyższe powtarzane przez $N_\text{gen}$ generacji.
\end{enumerate}

\paragraph{Hiperparametry początkowe.}
\begin{center}
\begin{tabular}{lr}
\toprule
Parametr & Wartość \\ \midrule
Rozmiar populacji & 60 \\
Liczba generacji & 100 \\
Prawdopodobieństwo krzyżowania & 0{,}7 \\
Prawdopodobieństwo mutacji & 0{,}2 \\
Mecze na ocenę fitness & 200 rozdań / przeciwnik \\
\bottomrule
\end{tabular}
\end{center}

\subsection{Funkcja celu}
Maksymalizujemy oczekiwaną wygraną na rozdanie:
\[
f(\theta) = \frac{1}{|\mathcal{O}|}\sum_{o \in \mathcal{O}}
\frac{1}{N}\sum_{i=1}^{N} u_1\!\left(\pi_\theta, o, \omega_i\right),
\]
gdzie $\pi_\theta$ to polityka określona przez wagi $\theta$,
$\mathcal{O}$ to pula przeciwników, $u_1$ to wypłata gracza
prowadzonego przez $\pi_\theta$, a $\omega_i$ to losowe rozdanie kart.

\subsection{Strategia referencyjna: równowaga Nasha}
Kuhn Poker posiada jednoparametrową rodzinę równowag Nasha
indeksowaną wartością $\alpha \in [0, 1/3]$. Dla $\alpha = 1/6$
strategia gracza pierwszego ma postać:
\begin{itemize}
    \item J: \emph{bet} z prawdopodobieństwem $\alpha$, w przeciwnym razie \emph{check}.
    \item Q: zawsze \emph{check}; w odpowiedzi na bet \emph{call} z prawdopodobieństwem $\alpha + 1/3$.
    \item K: \emph{bet} z prawdopodobieństwem $3\alpha$.
\end{itemize}
Wartość gry dla gracza działającego pierwszego wynosi
$-1/18 \approx -0{,}0556$, niezależnie od wyboru $\alpha$.

\end{document}
